\begin{table}[htbp]
\centering
\begin{threeparttable}
\caption{Alternative Policy Variants (APV) — counterfactual price effects}
\label{tab:v3_apv}
\small
\begin{tabular}{lrrrrr}
\toprule
Class & $\bar p_{S_1}$ & V0 (v3) & V1 (50\%) & V2 (no entry) & V3 (10\% pref.) \\
\midrule
non-pharma & 0.760 & +0.2436 & +0.0994 [40.8\%] & +0.3532 [145.0\%] & +0.0103 [4.2\%] \\
pharma & 0.649 & +0.3055 & +0.0885 [29.0\%] & +0.5897 [193.0\%] & +0.0032 [1.0\%] \\
\bottomrule
\end{tabular}
\begin{tablenotes}\footnotesize
\item Price $\Delta$ relative to observed $S_1$ (open, Pre pool).
V0 = full 100\% SME set-aside with endogenous Post pool (main text).
V1 = partial 50/50 mix. V2 = SME-only but Pre pool (isolates intensive
margin). V3 = open auction with 10\% price preference for SMEs.
Bracketed values are each variant's $\Delta$ as \% of V0's $\Delta$,
so V1 = 50\% means ``half the v3 effect''.
\end{tablenotes}
\end{threeparttable}
\end{table}
